\chapter{Extintores Portátiles y Sistemas Fijos de Extinción}
\label{sec.cap6:sistemas_extincion}

La supresión del fuego demanda la aplicación inmediata y calculada del agente extintor idóneo sobre el combustible en
combustión. La protección contra incendios se articula en dos niveles complementarios: los equipos portátiles de primera
intervención (matafuegos manuales y sobre ruedas) diseñados para combatir focos incipientes, y las instalaciones fijas
automáticas de extinción orientadas a contener y sofocar siniestros de mayor escala en infraestructuras industriales y
electrónicas.

\section{Extintores Portátiles}
\label{sec.cap6:extintores}

Los extintores portátiles constituyen la primera línea activa de defensa disponible en cualquier establecimiento. Su
diseño, selección, ensayo, dotación y mantenimiento se encuentran regulados en la República Argentina por el Decreto
Reglamentario 351/79 (Capítulo 18 y Anexo VII) y las normas IRAM 3517-1 (selección e instalación) e IRAM 3517-2
(control, mantenimiento y recarga).

\subsection{Clasificación según el Agente Extintor}

El agente almacenado en el cilindro define la eficacia química y la compatibilidad operativa con cada tipo de fuego:

\begin{itemize}
    \item Agua Presurizada y Agua Pulverizada: En su configuración de chorro pleno, actúan exclusivamente sobre fuegos
          Clase A mediante enfriamiento endotérmico. En su versión de agua desmineralizada con tobera pulverizadora especial,
          permite su uso seguro en recintos con tensión eléctrica moderada sin riesgo de electrocución ni corrosión por sales.
    \item Espuma Mecánica (AFFF y AR-AFFF): Solución acuosa con concentrados sintéticos que forman una
          película acuosa flotante. Apto para fuegos Clase A y Clase B. La película
          sella la superficie líquida, bloqueando la emisión de vapores inflamables y enfriando el
          combustible. 
    \item Polvo Químico Seco (PQS) Triclase ABC: Compuesto en un $75\percent$ a $90\percent$ por fosfato monoamónico
          ($\mathrm{NH_4H_2PO_4}$) finamente micronizado con agentes siliconados hidrofugantes. Ofrece una versatilidad
          sobresaliente: al ser proyectado sobre brasas sólidas (Clase A), el calor descompone la sal a partir de
          $180\Celsius$, tapona los poros y aísla el carbono del oxígeno; en líquidos inflamables (Clase B), el polvo
          detenie la reacción en cadena y por ser un polvo aislante, no conduce la electricidad (Clase C). No obstante
          es higroscópico y altamente corrosivo, por lo que su descarga es destructiva.
    \item Polvo Químico Seco BC: Formulado a base de bicarbonato de sodio ($\mathrm{NaHCO_3}$) o bicarbonato de potasio
          ($\mathrm{KHCO_3}$). Desarrolla una efectividad superior en fuegos de hidrocarburos y gases sin
          dejar una capa adherente vidriada, facilitando su limpieza en maquinarias pesadas.
    \item Dióxido de Carbono ($CO_2$): Gas comprimido a alta presión en cilindros de acero sin costura, encontrándose en
          fase mixta líquida y gaseosa a una presión de vapor de $50\text{ a }60\,\mathrm{bar}$ a $20\Celsius$. Al ser
          liberado a través de una tobera difusora cónica, experimenta una expansión brusca, proyectando gas a
          $-78.5\Celsius$. Extingue por desplazamiento del oxígeno y enfriamiento superficial. Al ser un gas inerte,
          limpio y rigurosamente dieléctrico, no deja ningún residuo químico ni humedad, siendo el agente idóneo para
          tableros eléctricos, consolas de mando, bancos de pruebas y salas de servidores.
    \item Agentes Limpios Halocarbonados: Sustancias químicas en recipientes presurizados. Se descargan en forma de líquido vaporizante de
          evaporación rápida que absorbe calor e interrumpe la combustión sin provocar el choque térmico extremo del $CO_2$
          ni dejar partículas conductoras, protegiendo instrumental óptico y electrónico sensible.
    \item Polvos Especiales para Metales Clase D: Compuestos a base de cloruro de sodio con aditivos poliméricos
          termoplásticos (agentes tipo MET-L-X) o polvos de grafito enriquecidos con sales de cobre (Lith-X). Se depositan
          suavemente sobre el metal fundido mediante palas o toberas de baja velocidad, sinterizándose para conformar una
          coraza asfixiante que impide el acceso de oxígeno atmosférico.
    \item Químico Húmedo Clase K: Solución acuosa alcalina de acetato y citrato de potasio. Al proyectarse en forma de niebla
          sobre aceites de cocina a más de $300\Celsius$, reacciona con los ácidos grasos libres formando un manto espeso de
          jabón no combustible que extingue el fuego y refrigera el fluido impidiendo su reignición espontánea.
\end{itemize}

\subsection{Clasificación según el Modo de Presurización}

En función del método mecánico empleado para impulsar el agente hacia el exterior, se dividen en dos grupos:

\begin{itemize}
    \item Presurización Permanente (Presión Contenida): El cilindro alberga conjuntamente el agente extintor y el gas
          propulsor (generalmente nitrógeno seco comercial) sometidos a presión constante en reposo. Incorporan un
          manómetro indicador frontal cuya aguja situada en la zona verde confirma la operatividad del sistema.
    \item Presurización No Permanente (Cápsula o Botellín Impulsor): El recipiente que contiene el agente se halla a presión
          atmosférica normal durante el almacenamiento. La presurización se produce únicamente en el instante del uso,
          cuando el operador presiona un gatillo percutor que perfora el diafragma de un pequeño cartucho de gas presurizado
          ($CO_2$ o nitrógeno) montado en el exterior o interior del extintor. Carecen de manómetro en el cuerpo principal.
\end{itemize}

\subsection{Clasificación según Capacidad y Movilidad}

\begin{itemize}
    \item Matafuegos Portátiles Manuales: Unidades con masas totales de carga de entre $1\kg$ y $10\kg$ (o hasta
          $10\unit{\liter}$ en equipos hídricos), diseñados para ser transportados y accionados manualmente por un solo
          operador.
    \item Matafuegos Rodantes (Sobre Ruedas o Carretillas): Unidades de mayor envergadura, con masas de carga de $25\kg$,
          $50\kg$ o hasta $100\kg$, montadas sobre un bastidor con ruedas y provistas de mangueras de descarga de hasta
          $10\m$, destinadas a la protección de áreas industriales de alto riesgo como subestaciones, parques de tanques y
          playas de carga y descarga de combustibles.
\end{itemize}

\subsection{Potencial Extintor y Dotación Reglamentaria}

El potencial extintor certifica numéricamente la capacidad de sofocación del matafuego en pruebas de
fuego normalizadas de laboratorio. Para la Clase A, los valores van de 1-A a 40-A, donde un extintor 2-A extingue el
doble de masa de madera que uno 1-A. Para la Clase B, el potencial varía de 1-B a 120-B, y para la Clase C, la letra se
consigna sin coeficiente numérico, acreditando únicamente la no conductividad dieléctrica del equipo y su tobera.

El Decreto Reglamentario 351/79 (Anexo VII) establece los siguientes requisitos de dotación en planta:
\begin{itemize}
    \item Para fuegos Clase A: La dotación mínima de unidades de extinción se computa en función de la superficie del
          sector y de la Carga de Fuego calculada según la ecuación \ref{eq.cap4:carga_fuego}. La distancia máxima de
          recorrido real desde cualquier punto del sector hasta alcanzar un extintor Clase A no puede exceder los $20\m$.
    \item Para fuegos Clase B: Se selecciona el potencial extintor individual conforme a la tabla de riesgo del decreto,
          exigiéndose una distancia máxima de recorrido no mayor a $15\m$.
    \item Ubicación y Altura de Montaje: Los extintores se ubicarán en lugares visibles, de fácil acceso y libres de
          obstáculos, sobre los caminos habituales de tránsito. La manija de sujeción debe instalarse a una altura
          comprendida entre $1.20\m$ y $1.50\m$ respecto al nivel de piso terminado. Cada puesto debe contar con su
          correspondiente chapa baliza reglamentaria de franjas diagonales rojas y blancas, y la indicación en tipografía
          clara del tipo de matafuego y las clases de fuego autorizadas.
\end{itemize}

\section{Sistemas Fijos de Extinción de Incendios}
\label{sec.cap6:sistemas_fijos}

Cuando las dimensiones del establecimiento o la concentración calórica superan la capacidad de extinción manual, el
Decreto Reglamentario 351/79 (Anexo VII, Inciso 4 y Cuadro de Protección y Servicio contra Incendios) y las normas IRAM
pertinentes (como IRAM 3597 para hidrantes, e IRAM 3501 e IRAM 3546 para certificación y mantenimiento) exigen la
implementación de sistemas fijos de extinción automáticos o manuales.

\subsection{Redes de Hidrantes e Instalaciones de Agua contra Incendios}

Las redes de hidrantes constituyen la columna vertebral hidráulica de la extinción industrial, reguladas técnicamente
en el país por la norma IRAM 3597 y las condiciones de servicio del Decreto 351/79. La instalación consta de:
\begin{itemize}
    \item Reserva de Agua Exclusiva: Cisterna dedicada con volumen suficiente para suministrar el caudal de diseño durante
          un lapso ininterrumpido de $60\text{ a }120\unit{\minute}$. La toma de agua de incendios se ubica en el fondo de
          la cisterna, por debajo de la reserva de agua de uso sanitario o industrial, asegurando su disponibilidad plena.
    \item Sala de Bombas: Integra una bomba principal impulsada por motor eléctrico, una motobomba de respaldo impulsada
          por motor diésel con arranque automático por banco de baterías (indispensable para operar ante cortes de red
          eléctrica), y una pequeña electrobomba presurizadora denominada bomba jockey. Esta última opera de manera
          intermitente para compensar microfugas en las cañerías y sostener la red presurizada permanentemente entre
          $8\,\mathrm{bar}$ y $10\,\mathrm{bar}$.
    \item Bocas de Incendio Equipadas (BIE): Gabinetes metálicos instalados en muros a distancias regulares, que contienen
          mangueras de fibras sintéticas de $38\mm$ ($1.5''$) o $63.5\mm$ ($2.5''$), acoples rápidos tipo Storz o
          mandrilados, llaves de ajuste y lanzas regulables de triple propósito (chorro pleno, niebla y cierre).
    \item Conexión Siamesa de Fachada: Colector perimetral exterior con válvulas de retención accesible desde la calzada,
          que permite a las autobombas de los bomberos inyectar agua presurizada directamente a la red del edificio.
\end{itemize}

\subsection{Sistemas de Rociadores Automáticos}

Los sistemas de rociadores automáticos operan de manera autónoma sin intervención humana. Se distribuyen sobre el cielorraso
en una red de cañerías que alimentan cabezas rociadoras individuales. Cada rociador se encuentra obturado por un bulbo de
vidrio de cuarzo sellado que contiene un líquido con alto coeficiente de dilatación térmica y una burbuja de aire. Al subir
la temperatura local a causa de un foco de incendio, el líquido se dilata hasta disolver la burbuja y reventar la ampolla,
liberando un orificio por donde el agua impacta contra un deflector estriado que genera un patrón de lluvia en paraguas.

La temperatura nominal de disparo del bulbo térmico responde a un código internacional de colores normalizado:
\begin{itemize}
    \item Naranja: $57\Celsius$
    \item Rojo: $68\Celsius$ (umbral estándar comercial y corporativo)
    \item Amarillo: $79\Celsius$
    \item Verde: $93\Celsius$
    \item Azul: $141\Celsius$
    \item Púrpura: $182\Celsius$
    \item Negro: Mayor o igual a $204\Celsius$
\end{itemize}

Según la modalidad de carga de las cañerías, los sistemas de rociadores se clasifican en:
\begin{itemize}
    \item Tubería Húmeda: Las cañerías permanecen continuamente llenas de agua presurizada. La rotura
          del bulbo produce la descarga inmediata de agua. Ofrece la máxima simplicidad y fiabilidad, pero no puede
          instalarse en recintos donde la temperatura descienda por debajo de $4\Celsius$ por riesgo de congelamiento y
          rotura de cañerías.
    \item Tubería Seca: Las cañerías están cargadas con aire comprimido o nitrógeno a presión. Una
          válvula diferencial retiene el agua en la sala de bombas. Al romperse un rociador, el escape del aire descomprime la
          línea, permitiendo que la válvula diferencial se abra y el agua inunde la red hasta salir por el rociador abierto.
          Apto para cámaras frigoríficas o plantas exteriores.
    \item Sistemas de Acción Previa: La cañería se encuentra vacía o bajo aire a baja presión de
          supervisión. La válvula principal de control requiere dos eventos independientes para autorizar la entrada de agua:
          la confirmación de un detector de humo de la central electrónica y la rotura térmica del bulbo. Esta doble
          condición de disparo resulta indispensable en centros de datos, cuartos de comunicaciones y salas eléctricas,
          evitando que un golpe accidental sobre un rociador provoque el anegamiento accidental de equipos sensibles.
    \item Sistemas de Diluvio: Emplean cabezas rociadoras abiertas en su totalidad (sin ampolla fusible).
          Una válvula automática de diluvio, comandada por un sistema de detección de llama o humo, inunda toda el área
          simultáneamente. Su aplicación se orienta a transformadores exteriores de subestaciones y áreas petroquímicas.
\end{itemize}

\subsection{Sistemas de Agua Nebulizada}

Emplean bombas de alta presión ($> 35\text{ a }140\,\mathrm{bar}$) y boquillas micronizadoras especiales para generar una
niebla densa compuesta por microgotas con diámetros inferiores a $100\,\mu\mathrm{m}$. Las microgotas presentan una relación
superficie-volumen extremadamente alta, provocando una tasa de absorción calórica colosal que enfría rápidamente la llama.
Simultáneamente, al alcanzar los $100\Celsius$, cada litro de agua líquida se expande unas $1700$ veces en forma de vapor,
desplazando localmente el comburente sin reducir el oxígeno en el resto del recinto. Requiere una fracción del volumen de
agua de los rociadores tradicionales, minimizando daños residuales en salas de generadores y turbinas.

\subsection{Sistemas de Inundación Total por Dióxido de Carbono}

Consisten en baterías centralizadas de cilindros de $CO_2$ conectados a una red de colectores y toberas de descarga
abiertas. Diseñados para descargar el agente en menos de $60\s$, alcanzando concentraciones volumétricas del $34\percent$ al
$75\percent$ en el recinto protegido.

Dado que una concentración de $CO_2$ en aire superior al $8\percent$ a $10\percent$ induce la pérdida súbita de conciencia y la
asfixia mortal en humanos, estos sistemas exigen estrictos enclavamientos de seguridad:
\begin{itemize}
    \item Válvula de bloqueo mecánico con candado de seguridad que debe cerrarse antes del ingreso
          de personal técnico de mantenimiento a la sala.
    \item Retardo neumático de descarga de $30\text{ a }60\s$, acoplado a sirenas neumáticas y ópticas de pre-descarga que
          ordenan la evacuación inmediata antes de liberar el gas.
    \item Enclavamiento automático con paro del sistema HVAC y cierre de persianas cortafuego para asegurar el confinamiento
          del gas y evitar su fuga a sectores adyacentes.
\end{itemize}

\subsection{Sistemas de Supresión mediante Agentes Limpios Gaseosos}

En centros de datos de alta disponibilidad, salas de servidores, salas limpias y centros de control industrial, la
continuidad operativa exige extinguir el fuego sin dañar los componentes electrónicos ni dejar residuos corrosivos. Para
estos entornos se emplean sistemas de agentes limpios gaseosos:

\begin{itemize}
    \item Gases Halocarbonados Sintéticos: Comprende el HFC-227ea (conocido comercialmente como FM-200) y la fluorocetona
          FK-5-1-12 (Novec 1230). Se descargan en lapsos de hasta $10\s$. Su mecanismo de extinción es físico: absorben la
          energía térmica a escala molecular, disminuyendo la temperatura de la llama por debajo del umbral de combustión.
          Son químicamente estables, dieléctricos y no degradan la capa de ozono. Las concentraciones de diseño se sitúan
          por debajo del nivel sin efecto adverso observable,
          garantizando márgenes de seguridad para las personas si una descarga se produce antes de concluir la evacuación.
    \item Gases Inertes Puros y Mezclas: Sistemas como el Inergen (IG-541, compuesto por $52\percent\,N_2$, $40\percent\,Ar$
          y $8\percent\,CO_2$) o el nitrógeno puro (IG-100). Sofocan el fuego reduciendo la concentración de oxígeno del
          $21\percent$ a un rango del $12\percent\text{--}14\percent$. La inclusión del $8\percent$ de $CO_2$ en la mezcla
          IG-541 estimula de forma refleja los quimiorreceptores del centro respiratorio humano, incrementando el volumen
          minuto respiratorio y permitiendo que el cerebro reciba el oxígeno indispensable durante el tiempo de escape.
\end{itemize}

\subsection{Protección Electrónica Avanzada: Dispositivos de Detección de Falla de Arco (AFDD)}

Complementando a los sistemas de extinción tradicionales, la norma IEC 62606 estipula el empleo de Dispositivos de
Detección de Defectos de Arco (AFDD). Estos módulos electrónicos de tablero integran microcontroladores dedicados que
muestrean continuamente la forma de onda de la corriente eléctrica a alta frecuencia. Mediante análisis espectral y
transformadas rápidas de Fourier (FFT), reconocen las firmas de ruido estocástico y las asimetrías de cruce por cero
características de las descargas de plasma en arcos en serie. Al detectar la firma de falla de arco, el dispositivo ordena
la apertura del interruptor en tiempos del orden de los milisegundos, suprimiendo la fuente de calor antes de que la
temperatura alcance el punto de autoignición de los plásticos.

